% Generated from validated Article Draft Result. Do not hand-edit; revise the structured draft instead.
\noindent\textbf{vLLM、SGLang、FlashInferのreleaseが続く一方、研究側ではKV cache再利用やMoE expert reuseが掘り下げられた。モデルの外側にあるserving stackは、単一kernelではなく複数layerの協調問題になっている。}\autocite{src-3707e05fd2a3,src-5a693cfc6852,src-69c5523286bd,src-74a7052acd7a,src-77e29e3fac8f,src-80b8d67933c0,src-ba75040df3fc,src-c97320147e8c}

\subsection{5月の変化を一つのsystemとして読む}

同じモデルでも、どのruntime・attention backend・cache方針・parallelism・kernelで動かすかによって実運用の制約は大きく変わる。5月はvLLM、SGLang、FlashInferがそれぞれ継続的に更新され、研究側でもlong-context KV reuseやmemory-constrained MoE inferenceが別の角度から扱われた。\autocite{src-3707e05fd2a3,src-5a693cfc6852,src-69c5523286bd,src-74a7052acd7a,src-77e29e3fac8f,src-80b8d67933c0,src-ba75040df3fc,src-c97320147e8c}


\subsection{vllm-project/vllm May 2026 release series}

vLLMは5月4日のv0.20.1、15日のv0.21.0、29日のv0.22.0と更新を重ねた。release notesではDeepSeek V4 hardening、KV/offload、attention backend、Model Runner V2、speculative/disaggregated servingなどが扱われており、serving engineが複数のmodel familyとruntime techniqueを同時に吸収している。\autocite{src-77e29e3fac8f,src-80b8d67933c0,src-c97320147e8c}


\subsection{sgl-project/sglang May 2026 release series}

SGLangは5月5日のv0.5.11と16日のv0.5.12で、speculative decoding、disaggregated prefix caching、model support、DeepSeek V4 inference pathなどを更新した。個別のthroughput値ではなく、prefix/cacheとmodel-specific executionが同じrelease seriesで進んでいる点に注目したい。\autocite{src-5a693cfc6852,src-69c5523286bd}


\subsection{Release v0.6.12}

FlashInfer v0.6.12は5月29日に公開され、release notesにはNVFP4 quantization、MoE kernel、GQA/MLA path、model/backend supportが並ぶ。kernel layerでもquantization・attention・MoEが密接に結び付いており、frontier model対応が低level runtimeへ波及している。\autocite{src-74a7052acd7a}


\subsection{Adaptive KV Cache Reuse for Fast Long-Context LLM Serving}

CacheTuneは、厳密なprefix一致だけに頼らず、attention関係を保つため必要なcontextを適応的に再計算しながらKV cacheを再利用する設計を提案する。これは著者による研究結果であり、一般的なspeedup保証ではないが、long-context servingで「何を再利用できるか」自体を最適化対象にする方向を示す。\autocite{src-3707e05fd2a3}


\subsection{ReMoE: Boosting Expert Reuse through Router Fine-Tuning in Memory-Constrained MoE LLM Inference}

ReMoEはmemory-constrainedなMoE inferenceでexpert reuseを高めるためrouter fine-tuningを利用する。MoE servingでは計算量だけでなくexpert weightのresident／transfer costが問題になるため、model architectureとruntime memory制約の境界を跨ぐ研究として位置付けられる。\autocite{src-ba75040df3fc}


\subsection{Theme Synthesis}

ここで重要なのは、個々のproject benchmarkを一つの勝敗表へ混ぜないことだ。release noteは実装変更の一次資料だが性能値はproject条件に依存し、論文のspeedupも著者が選んだmodel・hardware・workload・baselineに依存する。5月の意味は、最適化対象がcacheからexpert、kernel、runtimeまでstack全体へ広がった点にある。\autocite{src-3707e05fd2a3,src-5a693cfc6852,src-69c5523286bd,src-74a7052acd7a,src-77e29e3fac8f,src-80b8d67933c0,src-ba75040df3fc,src-c97320147e8c}


\begin{claimboundary}
Claim Boundary: GitHub release notesはrelease chronologyと記載された実装変更を確認する一次資料だが、projectが示すperformanceやcross-hardwareな一般化は本稿で独立再現していない。\autocite{src-3707e05fd2a3,src-5a693cfc6852,src-69c5523286bd,src-74a7052acd7a,src-77e29e3fac8f,src-80b8d67933c0,src-ba75040df3fc,src-c97320147e8c}
\end{claimboundary}

\begin{claimboundary}
Claim Boundary: 原論文は研究内容を確認する一次資料だが、benchmark結果・speedup・attack success・quality比較は著者報告であり、本稿で独立再現した結果ではない。\autocite{src-3707e05fd2a3,src-5a693cfc6852,src-69c5523286bd,src-74a7052acd7a,src-77e29e3fac8f,src-80b8d67933c0,src-ba75040df3fc,src-c97320147e8c}
\end{claimboundary}
